\chapter{Обзор предметной области и постановка задачи}

\section{Постановка задачи траекторного управления}

В работе рассматривается движение квадрокоптера вдоль заданной пространственной кривой. Требуется обеспечить не просто выход объекта на траекторию, а устойчивое слежение за ней с допустимыми ошибками и с контролируемой скоростью продвижения вдоль кривой. Иначе говоря, объект должен двигаться по геометрически заданному маршруту так, чтобы поперечное отклонение не выходило за допустимые пределы, а динамика оставалась корректной на всем интервале моделирования.

Если зафиксировать параметрическую скорость движения вдоль кривой, задача сводится к построению устойчивого регулятора слежения. Такой подход удобен и хорошо согласуется с классической схемой синтеза управления. Однако на практике он оказывается излишне жестким. Для прямолинейных и слабо искривленных участков траектории фиксированная скорость часто выбирается с большим запасом, тогда как на участках с высокой кривизной тот же самый режим уже может приводить к росту ошибки. Поэтому в рамках ВКР задача ставится шире: базовый регулятор должен обеспечивать слежение, а дополнительный интеллектуальный модуль должен подбирать значение $V^{\ast}$ в зависимости от текущего состояния системы.

\section{Базовый подход к решению}

В качестве основы используется алгоритм согласованного управления по выходу, описанный в работе С.\,А.~Ким~\cite{kim_algorithms}. Его ключевая идея состоит в том, что управление строится не напрямую по координатам объекта, а по регулируемым ошибкам, заданным относительно ближайшей точки опорной кривой. Для этого вводится наблюдатель ближайшей точки, а сами ошибки формируются в сопутствующей системе координат, связанной с геометрией траектории.

Такой подход удобен по нескольким причинам. Во-первых, он естественно связывает динамику объекта с геометрией траектории. Во-вторых, позволяет отдельно анализировать продольную и поперечную составляющие ошибки. В-третьих, дает понятную структуру закона управления, которая хорошо переносится в программную реализацию и подходит для дальнейшего расширения.

В проекте этот базовый алгоритм не заменяется, а используется как устойчивый нижний уровень управления. Интеллектуальный модуль работает поверх него и влияет только на выбор целевой скорости продвижения вдоль кривой.

\section{Объект управления и принятые ограничения}

Объектом управления выбран квадрокоптер, движущийся в трехмерном пространстве. Такой выбор удобен по двум причинам: с одной стороны, модель достаточно содержательная и отражает особенности нелинейной динамики, с другой стороны, она остается достаточно компактной для проведения серий численных экспериментов.

В работе используется расширенная 16-мерная модель состояния
\begin{gather}
    x =
    \begin{bmatrix}
        p_x & p_y & p_z &
        v_x & v_y & v_z &
        \phi & \theta & \psi &
        \dot{\phi} & \dot{\theta} & \dot{\psi} &
        \bar{u}_1 & \rho_1 & u_2 & \rho_2
    \end{bmatrix}^{\top},
\end{gather}
где первые компоненты описывают положение, линейную скорость и ориентацию объекта, а последние четыре компоненты вводятся как интеграторные расширения, необходимые для реализации регулятора.

Заданная траектория рассматривается в параметрическом виде
\begin{gather}
    p^{\ast}(s) = \begin{bmatrix} x^{\ast}(s) & y^{\ast}(s) & z^{\ast}(s) \end{bmatrix}^{\top}.
\end{gather}
В текущей реализации используются кривые, для которых норма касательного вектора постоянна:
\begin{gather}
    \left\|\frac{d p^{\ast}(s)}{d s}\right\| = \mathrm{const}.
\end{gather}
Это ограничение связано не с удобством записи, а с тем, что именно для такого класса траекторий корректно работает выбранный базовый алгоритм.

В качестве тестовых и обучающих траекторий используются три семейства кривых:
\begin{enumerate}
    \item пространственная прямая;
    \item горизонтальная окружность;
    \item спираль.
\end{enumerate}
Этого набора достаточно, чтобы покрыть как простой случай движения без выраженной кривизны, так и более сложные сценарии, в которых допустимая скорость зависит от геометрии маршрута.

\section{Регулируемые переменные и критерии качества}

Для описания качества слежения используются ошибки, связанные с положением объекта относительно ближайшей точки траектории. В текущей реализации ключевыми величинами являются продольная ошибка, поперечная ошибка и ошибка по курсу. На их основе формируется вектор регулируемых переменных
\begin{gather}
    \lambda_1 =
    \begin{bmatrix}
        s_{\mathrm{arc}} - s_{\mathrm{ref}} \\
        e_1 \\
        e_2 \\
        \delta_{\phi}
    \end{bmatrix},
\end{gather}
где $s_{\mathrm{arc}}$ соответствует фактически пройденной длине дуги вдоль траектории, а $s_{\mathrm{ref}}$ задает опорное продвижение, формируемое по целевой скорости $V^{\ast}$.

С практической точки зрения основными критериями качества являются:
\begin{enumerate}
    \item сходимость поперечной ошибки $e_2$ к нулю;
    \item ограниченность продольной ошибки и ошибки синхронизации;
    \item отсутствие численной неустойчивости при моделировании;
    \item возможность увеличивать скорость движения без выхода за допустимые ошибки.
\end{enumerate}

Именно последний пункт определяет логику всей дальнейшей работы. Базовый регулятор обеспечивает устойчивость, а интеллектуальная часть должна ответить на вопрос, какое значение $V^{\ast}$ допустимо в текущих условиях.

\section{Выводы по главе}

Таким образом, задача ВКР сводится к построению двухуровневой схемы управления. На нижнем уровне находится аналитический регулятор траекторного слежения, использующий математическую модель объекта и геометрию опорной кривой. На верхнем уровне располагается модуль адаптивного выбора параметрической скорости, который должен повышать эффективность движения, не разрушая устойчивость базового контура. В следующих главах рассматриваются программная реализация этой схемы и результаты численного исследования.
